\subsection{The solvable quotient obtained from the raw orthogonal blocks}

\begin{statement}{D23.F37}{Proved}{The raw block quotient restores the cubic Burau scalar}
\end{statement}
In the raw degree-$8064$ action of D23.H6, lift the twelve
projective blocks of D23.F30. Each lifted block has $672$ points.
The permutation image $H$ of the stabilizer of one block has
\[
 |H|=39191040,\qquad |H'|=|H''|=6531840,\qquad H/H'\cong C_6.
\]
The induced action on twelve copies of the regular coset set
$H/H'$ is a degree-$72$ quotient $Q_{72}$ of the raw braid image.
Its derived series has orders
\[
              331776,55296,13824,6912,768,1.
\]
It is exactly, with the original three braid generators marked,
the literal matrix group generated by FP.36.2 over
$R=\mathbb Z/4[\omega]/(\omega^2+\omega+1)$, before quotienting
by scalar matrices. The natural map to D23.F31's signed
$24$-point quotient has central kernel
\[
                      \langle\omega I\rangle\cong C_3,
\]
generated by the four-strand full twist $(ABC)^4$.
This kernel is contained in the derived subgroup: the two derived
orders are $55296$ and $18432$, respectively. In particular it
would be incorrect to describe this extension as adding only an
abelian character of the whole braid image.

Nevertheless every one of the five marked evaluation pair images
maps isomorphically to its signed degree-$24$ counterpart:
\[
\begin{array}{c|rrrrr}
\text{pair}&12&23&34&45&51\\\hline
\text{image order}&192&192&144&48&144.
\end{array}
\]
Thus no value at any one of these five pairs gains an independent
scalar choice in the raw solvable quotient. The marked cubic
$216$ quotient still does not factor through its first-pair image,
whereas the marked harmonic $S_4$ quotient does factor through
its $34$ image. This concerns the induced solvable quotient only,
not all information in the raw degree-$8064$ action.

\begin{proof}
Write $q$ for the block permutation of a raw braid permutation
$g$, and choose twelve raw braid transversals $T_i$ from block
zero to block $i$. The certificate obtains them by an exhaustive
traversal of the twelve-block action. The Schreier elements
\[
                       T_{q(i)}^{-1}gT_i
\]
for the three generators and all twelve blocks generate the
stabilizer. Restricting these elements to the original $672$-point
block gives thirteen distinct nonidentity generators of its full
image $H$. The exact permutation-group order and two successive
derived orders are as displayed. Equality $H'=H''$ proves
perfectness of $H'$. The quotient has order six and is abelian,
so it is cyclic.

Let $L(g,i)$ denote the restriction of
$T_{q(i)}^{-1}gT_i$ to the base block. On pairs $(i,hH')$ define
\[
                  g(i,hH')=(q(i),L(g,i)hH').
\]
The identity
$L(gg',i)=L(g,q'(i))L(g',i)$ proves the action law.
Since $H'$ is normal in $H$, this is well defined. Enumerating
the six cosets by exact membership in $H'$ therefore constructs
an actual quotient of the raw image, not an independently chosen
finite representation. The three resulting degree-$72$
permutations generate a group of the stated order and complete
derived series.

For the matrix identification, encode $a+b\omega\in R$ by
$a+4b$, $0\leq a,b<4$, and use the literal FP.36.2 matrices
\[
 A=\begin{pmatrix}1-\omega&\omega&0\\1&0&0\\0&0&1\end{pmatrix},
 \quad B=\begin{pmatrix}1&0&0\\0&1-\omega&\omega\\0&1&0\end{pmatrix},
 \quad C=\begin{pmatrix}1&0&-1\\0&1&-1\\0&0&-\omega\end{pmatrix}.
\]
Traverse the full degree-$72$ permutation group while multiplying
these three matrices along the identical generator edges. There
are $331776$ distinct permutation elements. Every one of the
$995328$ generator edges gives the same matrix whenever a
permutation is reached again, and all $331776$ resulting
matrices are distinct. This gives a well-defined bijective marked
homomorphism onto the literal matrix group. It is an exact
relation check, not recognition by matching group orders.
The full twist maps to $\omega I$.

D23.F31 identifies the signed quotient with the same matrix group
modulo $\langle\omega I\rangle$. Since both identifications use
the same three generators, their diagram commutes. Hence the
projection has precisely the displayed central kernel. The
certificate independently verifies this projection by joining
the two marked permutation triples: the joint image has order
$331776$, equal to $|Q_{72}|$. It also computes the two derived
orders. Their quotient is three, which proves that the entire
kernel lies in $Q_{72}'$.

Finally form each of the five literal pairs
$X_1=A^2$, $X_2=B^2$, $X_3=C^2$, $X_4=(ABA)^2$,
$X_5=(CBC)^2$ in both quotients. For every pair the two individual
orders and the order of their marked joint image are equal to
the corresponding entry in the table. Both projections are onto,
so equality proves that each pair projection is an isomorphism.
In particular the same word at that pair has a unique raw
solvable value above its signed value.

For the final marked quotient tests, join the $12$ pair with the
explicit nine-point affine $216$ detector of D23.P17. Its joint
order is $3456$, strictly larger than $192$, so that marked map
does not factor through the first-pair image. Joining the $34$
pair with the literal $S_4$ pair $((123),(1234))$ leaves order
$144$, so its marked map does factor. These are tests of the
specified two-generator maps, not statements inferred only from
the abstract orders of potential target groups.
\end{proof}

\paragraph{Exact certificates and limitation.}
The construction and derived groups are in
\path{agents/five_level_six/hurwitz_raw_block_solvable.py};
the marked quotient and all five pair checks are in
\path{hurwitz_raw_solvable_markings.py}; the complete matrix
isomorphism replay is \path{hurwitz_raw_burau.py}.
Each has its JSON record, and the block construction also retains
all transversals, cosets, local transitions and generators in NPZ.
No abstract identification of $H$ from its order is required,
and no assertion that this is the maximal solvable quotient of
the entire raw braid image is made. Raw perfect-kernel and
global common-word constraints remain available beyond this
particular quotient.
